%%% important comment keywords %%%
% CAVEAT: -> a warning concerning some compilation issues that may break your system in some situation
% ADJUST: -> manual tweaks that you may want to change
% TODO: -> to-do list
% HACK: -> Savy macros to easily use a given packages

%%%%%%%%% bibliography %%%%%%%%%
\usepackage[portuguese]{babel}
\usepackage[style=numeric-comp,backend=biber]{biblatex} % tells LaTeX to use the biber tool instead of the traditional bibtex tool for sorting and formatting your bibliography. Biber is a more advanced and powerful tool than bibtex and it provides more features, such as advanced sorting, Unicode support and support for various data sources.
\addbibresource{refs.bib} % add reference file, in the document, add \printbibliography to print bibliography
\addbibresource{refs_additional.bib}

%%%%%%%%% references (cite and glossary) %%%%%%%%%
\makeatletter
% some classes (e.g., beamer) loads hyperref by default. So it is sensible to check whether it wasn't already loaded to avoind clashes
\@ifpackageloaded{hyperref}{
    % The package is already loaded
}{
    % The package is not loaded, so load it now
    \usepackage[colorlinks=true,hyperindex=false]{hyperref} % make references blue and clickable
}
\makeatother
\hypersetup{ % you can use allcolors=blue in hyperref options instead this setup
    linkcolor=black,          % For \gls{} entries
    citecolor=blue,          % For \cite{} entries
}
\usepackage{glossaries} % \newglossaryentry{} entries. ADJUST: If you want to make the glossary entries nonclickable, load this package before hyperref
% CAVEAT: don't use `glossaries-extra`, otherwise the conditional glossary entry will not work! % SEE: https://tex.stackexchange.com/questions/98494/glossaries-dont-print-single-occurences/230664#230664
\glsenableentrycount % enable \cgls, \cglspl, \cGls, \cGlspl
% NOTE: use \cgls-based instead of \gls-based macros as the former puts the initials iff. it was used more than once.
% HACK:
% \glsentrylong{pll} -> long form

% add hyphened entries
% SEE: https://tex.stackexchange.com/a/311882/211183
\glsaddkey
    {hyphenated}        % new key
    {\relax}            % default value if "hyphenated" isn't used in \newglossaryentry
    {\cglsentryhyph}     % analogous to \cglsentrytext
    {\cGlsentryhyph}     % analogous to \cGlsentrytext
    {\cglshyph}          % analogous to \cglstext
    {\cGlshyph}          % analogous to \cGlstext
    {\cGLShyph}          % analogous to \cGLStext

%%%%%%%%% list of authors %%%%%%%%%
\usepackage{authblk} % CAVEAT: this pkg will lead to compilation errors if you use it along with `\IEEEauthorblockN`, which is a command from the `IEEEtran` document class

%%%%%%%%% items, tables and figs %%%%%%%%%
\makeatletter
% only load enumitem if not already present (sbrt.cls provides \labelindent)
\@ifpackageloaded{enumitem}{}{\@ifundefined{labelindent}{\usepackage{enumitem}}{}}
\makeatother
\usepackage{xltabular}
\usepackage{float} % use the [H] option
\usepackage{tikz}
\usepackage{graphicx} % extends the basic LaTeX graphic capabilities: placement, scaling, rotation, file formats
\usepackage{subcaption} % provides a means of using the subfigures environment, allowing a figure to be subdivided into smaller parts
\usepackage{standalone} % place tikz environments or other material in own source files

%%%%%%%%% new features %%%%%%%%%
\usepackage{xparse} % create commands with optional arguments or complex argument structures
\usepackage{xstring} % testing a string's contents, extracting substrings, substitution of substrings

%%%%%%%%% math %%%%%%%%%
\newcommand\bmmax{0} \newcommand\hmmax{0} % fix the "Too many math alphabets used in version normal" error (see https://tex.stackexchange.com/a/243541/211183)
\usepackage{amsmath}
\usepackage{mathtools} % is an extension package to amsmath to use \DeclarePairedDelimiter
\usepackage{amsfonts}
\usepackage{amssymb}
\usepackage{esint} % provides extended integral symbols, giving you access to a wider variety of integral signs, such as double, triple, quadruple integrals
% HACK:
% \oiint -> contour integral
\usepackage{newtxmath} % for Greek variants (bold, nonitalic, etc...). E.g., `\boldsymbol{\betaup}`
\usepackage{stmaryrd} % provides new symbols, such as \llbracket \rrbracket
\usepackage{bm} % for writing tensor, e.g., $\bm{\mathcal{A}}$
% only load IEEEtrantools if the class did not already provide IEEEeqnarray
\makeatletter
\@ifpackageloaded{IEEEtrantools}{}{\@ifundefined{IEEEeqnarray}{\usepackage{IEEEtrantools}}{}}
\makeatother
\usepackage{siunitx} % typesetting units and unitless numbers SI (Système International d’Unités) units
% CAVEAT: you cannot use siunitx for nonnumerical values, e.g. \qty{\pi}{\meter\per\second}
% HACK:
% \qty -> number and unit, e.g., \qty{10}{\meter\per\second}
% \unit{} -> unit only, e.g., \unit{\meter\per\second}
% Examples:
% 1. \qtyrange[exponent-to-prefix=true, range-independent-prefix=true]{138e6}{2.9e9}{\hertz} -> 138 MHz to 2.9 GHz
% use, possible options:
% 1. "range-phrase=-" ->  change the word "to" to "--"
% 2. range-units=single -> chage "1m to 10m" to "1 to 10m"
% 3. mode=text % writes the unit in full
% 4. quantity-product={-} to create a compound adjective with quantities, e.g. "the 50-Hz signal..."
% use the [per-mode=symbol] option in the \unit or \qty command to change the writing of \per, or change the line below to set it globally
\sisetup{per-mode=symbol} % print the symbol "/" in "25/100"
\sisetup{tight-spacing=true} % makes the spacing in "5 x 10²" closer
% \squared -> ^2

\usepackage{physics} % several math macros, such as
\AtBeginDocument{\RenewCommandCopy\qty\SI} % CAVEAT: avoid incompatibilities between physics and siunitx packages when using the \qty{}{} command. The conflict happens because it is defined in both packages. Although physics has \SI{}{} as an alternative, it is deprecated and therefore you should use \unit{} and \qty{}{} instead of \si{} and \SI{}{}, respectively. See here (https://tex.stackexchange.com/questions/628183/how-to-avoid-qty-conflict-with-physics-and-siunitx) for more info.
% HACK:
% \tr -> trace
% \rank -> rank
% \expval -> angle brackets to inner product,〈a,b〉, or ensemble average〈s(t)〉
% \abs -> |x|
% \norm -> ||x||
% \eval -> evaluated bar x|_y=a
% \order -> Big-O notation
% \Re -> real
% \Im -> imaginary
% \dd[]{} -> differential
% \dv -> derivative
% \pdv -> partial derivative

%%%%%%%%% operators %%%%%%%%%
\let\oldemptyset\emptyset % change empty set
\let\emptyset\varnothing
% circular convolution a lá Oppenheim (if possible, prefer \circledast)
\newcommand*\circconv[1]{%
  \begin{tikzpicture}
    \node[draw,circle,inner sep=1pt] {#1};
  \end{tikzpicture}}
\newcommand{\intersection}{\bigcap\limits} % intersection operator

%%%%%%%%% delimiters %%%%%%%%%
\makeatletter % changes the catcode of @ to 11
\DeclarePairedDelimiter\ceil{\lceil}{\rceil} % ⌈x⌉
\let\oldceil\ceil
\def\ceil{\@ifstar{\oldceil}{\oldceil*}} % swap the asterisk and the nonasterisk behaviors

\DeclarePairedDelimiter\floor{\lfloor}{\rfloor} % ⌊x⌋
\let\oldfloor\floor
\def\floor{\@ifstar{\oldfloor}{\oldfloor*}}
\makeatother % changes the catcode of @ back to 12

%%%%%%%%% basic functions %%%%%%%%%
\newcommand{\adj}[1]{\ensuremath{\operatorname{adj}\left(#1\right)}} % adjugate matrix
\newcommand{\cof}[1]{\ensuremath{\operatorname{cof}\left(#1\right)}} % cofactor matrix
\newcommand{\eig}[1]{\ensuremath{\operatorname{eig}\left(#1\right)}} % eigenvalues
\newcommand{\nullspace}[1]{\ensuremath{\operatorname{N}\left(#1\right)}} % nullspace or kernel of the matrix
\newcommand{\nullity}[1]{\ensuremath{\operatorname{nullity}\left(#1\right)}} % nullity=dim(N(A))
\newcommand{\spn}[1]{\ensuremath{\operatorname{span}\left\{#1\right\}}} % span of a set of vectors, \span is a reserved word
\newcommand{\range}[1]{\ensuremath{\operatorname{C}\left(#1\right)}} % range or columnspace of a matrix=span(a1,a2, ..., an), where ai is the ith column vector of the matrix A
\newcommand{\unvec}[1]{\ensuremath{\operatorname{unvec}\left(#1\right)}} % unvectorize operator
\newcommand{\diag}[1]{\ensuremath{\operatorname{diag}\left(#1\right)}} % diagonal operator
\newcommand{\dom}[1]{\ensuremath{\operatorname{dom}\left(#1\right)}} % domain of the function
\newcommand{\frob}[1]{\ensuremath{\norm{#1}_\textrm{F}}} % Frobenius norm
\renewcommand{\dim}[1]{\ensuremath{\operatorname{dim}\left(#1\right)}} % dimension of a set

%%%%%%%%% customizable functions %%%%%%%%%
% arg min
\NewDocumentCommand{\argmin}{ e{_} m }{% \argmin_{x ∈ A}{f(x)} or \argmin{f(x)}
    \ensuremath{%
        \IfValueTF{#1}{%
            \mathop{\operatorname{arg\,min}}\limits_{#1}\;#2
        }{% else
            \operatorname{arg\,min}\;#2
        }
    }
}
% arg max
\NewDocumentCommand{\argmax}{ e{_} m }{% \argmax_{x ∈ A}{f(x)} or \argmax{f(x)}
    \ensuremath{%
        \IfValueTF{#1}{%
        \mathop{\operatorname{arg\,max}}\limits_{#1}\;#2
        }{% else
            \operatorname{arg\,max}\;#2
        }
    }
}
% limit in the mean (see Van Tress, chapter 6)
\NewDocumentCommand{\limean}{ E{_}{{}} m }{ % \limean_{x ∈ A}{f(x)} or \limean{f(x)}
    \ensuremath{
        \IfValueTF{#1}{%
            \underset{#1}{\operatorname{l.i.m.}\;#2}
        }{% else
            \operatorname{l.i.m.}\;#2
        }
    }
}
% mean
\NewDocumentCommand{\E}{ E{_}{{}} o m }{% \E{A}, \E_{u}{A}. If you want only the left or right part of \E (useful for breaking line) you can use \E_{u}[left-only]{A} or \E_{u}[right-only]{A}, respectively
    \ensuremath{%
        \IfValueTF{#2}{%
            \IfStrEq{#2}{left-only}{%
              \operatorname{E}_{#1}\left[#3\right.%
            }{% right-only
              \left.#3\right]%
            }%
          }{%
            \operatorname{E}_{#1}\left[#3\right]%
          }%
    }
}
% covariance
\NewDocumentCommand{\cov}{O{} m}{\ensuremath{\operatorname{cov}_{#1}\left[#2\right]}} % \cov[u]{x} or \cov{x}
% variance
\let\var\undefined\NewDocumentCommand{\var}{O{} m}{\ensuremath{\operatorname{var}_{#1}\left[#2\right]}} % \var[u]{x} or \var{x}
% vectorize operator
\let\vec\undefined\NewDocumentCommand{\vec}{O{} m}{\ensuremath{\operatorname{vec}_{#1}\left[#2\right]}} % \vec[u]{\mathbf{A}}, [u] is optional
% normalize all trigonometry functions (new and predefined) to put a parenthesis around the input argument
\RenewDocumentCommand{\sin}{s m}{% \sin{x} -> sin(x) \sin*{x} -> sin x
    \ensuremath{
    \IfBooleanTF{#1}{\operatorname{sin} #2}{\operatorname{sin}\left(#2\right)}%
    }
}
\RenewDocumentCommand{\cos}{s m}{% \cos{x} -> cos(x) \cos*{x} -> cos x
    \ensuremath{
    \IfBooleanTF{#1}{\operatorname{cos} #2}{\operatorname{cos}\left(#2\right)}%
    }
}
\RenewDocumentCommand{\tan}{s m}{% \tan{x} -> tan(x) \tan*{x} -> tan x
    \ensuremath{
    \IfBooleanTF{#1}{\operatorname{tan} #2}{\operatorname{tan}\left(#2\right)}%
    }
}
\NewDocumentCommand{\sinc}{s m}{% \sinc{x} -> sinc(x) \sinc*{x} -> sinc x
    \ensuremath{
    \IfBooleanTF{#1}{\operatorname{sinc} #2}{\operatorname{sinc}\left(#2\right)}%
    }
}
\NewDocumentCommand{\sgn}{s m}{% \sgn{x} -> sgn(x) \sng*{x} -> sgn x
    \ensuremath{
    \IfBooleanTF{#1}{\operatorname{sng} #2}{\operatorname{sgn}\left(#2\right)}%
    }
}
\RenewDocumentCommand{\sec}{s m}{% \sec{x} -> sec(x) \sec*{x} -> sec x
    \ensuremath{
    \IfBooleanTF{#1}{\operatorname{sec} #2}{\operatorname{sec}\left(#2\right)}%
    }
}

%%%%%%%%% minor adjusts in some mathematical symbols %%%%%%%%%
\DeclareMathAlphabet{\mathcal}{OMS}{cmsy}{m}{n} % newtxmath changes the font style of \mathcal. It prevents \mathcal from being changed

\let\oldforall\forall % put some spaces between the \forall command
\renewcommand{\forall}{\;\oldforall\;}

% fix the problem in putting integral limits above and below the int symbol
% SEE: https://tex.stackexchange.com/questions/148520/to-have-two-limits-in-double-integral
% HACK: \iint \limits_{-\infty}^{+\infty}
\makeatletter
%% make esint definition in line with amsmath
\@for\next:={int,iint,iiint,iiiint,dotsint,oint,oiint,sqint,sqiint,
  ointctrclockwise,ointclockwise,varointclockwise,varointctrclockwise,
  fint,varoiint,landupint,landdownint}\do{%
    \expandafter\edef\csname\next\endcsname{%
      \noexpand\DOTSI
      \expandafter\noexpand\csname\next op\endcsname
      \noexpand\ilimits@
    }%
  }
\makeatother

%%%%%%%%% comments, corrections, or revisions %%%%%%%%%
\usepackage{ifluatex} % handle lualatex-specific packages
\ifluatex
    \usepackage{luacolor} % colour support based on LuaTEX’s
    \usepackage[soul]{lua-ul} % provide underlining, strikethrough, and highlighting using features in LuaLATEX which avoid the restrictions imposed by other methods
\fi
\newcommand{\obs}[1]{\textcolor{red}{(#1)}} % observation
\newcommand{\secret}[1]{\makebox[0cm]{}} % inline comment
\newcommand{\sizecorr}[1]{\makebox[0cm]{\phantom{$\displaystyle #1$}}} % Used to seize the height of equation
\interfootnotelinepenalty=10000 % restrict the footnote to one page